% Source: Weinberg GR, physical PDF pp. 493--494
%         (printed pp. 469--470; p. 494 only through the text preceding
%         Section 15.1).
% Coverage: Chapter 15 opening material before Section 15.1.
% Figures/tables/footnotes: chapter epigraph; no figures, tables, or footnotes.
% Status: source-reviewed and compile-clean.
% Uncertainties: none.

\begin{flushright}
\begin{minipage}{0.46\textwidth}
\itshape
``Now entertain conjecture\\
of a time, when creeping\\
murmur and the poring\\
dark fills the wide vessel of\\
the universe.''\\[0.35em]
\normalfont\itshape William Shakespeare, \emph{King Henry\\
the Fifth}
\end{minipage}
\end{flushright}

\bigskip

In the last chapter we laid out the coordinates for a map of the universe in
space and time. Now we must begin to fill in this map with the islands of matter
and the seas of radiation that make up the physical contents of the universe.

For the most part, we shall continue to base our discussion on the assumption
of isotropy and homogeneity, now supplemented with Einstein's field equations.
The future of the universe then depends critically on its curvature: If the
universe is open, it will go on expanding forever, whereas if it is closed, its
present expansion will eventually cease and be succeeded by a general
contraction. The curvature in turn depends critically on the present energy
density \(\rho_0\); the universe is open or closed according to whether
\(\rho_0\) is less or greater than a critical value \(\rho_c\), of order
\(10^{-29}\,\mathrm{g/cm^3}\). It appears that \(\rho_0\) mostly arises from
the rest mass of ordinary matter---neutrons and protons. In this case, the
universe is open and \(\rho_0\) is less than \(\rho_c\) if the deceleration
parameter \(q_0\) is less than \(\tfrac12\), whereas the universe is closed and
\(\rho_0\) is greater than \(\rho_c\) if \(q_0\) is greater than
\(\tfrac12\); this justifies the emphasis on the measurement of \(q_0\)
throughout the last chapter. However, the observation that \(q_0\approx1\)
conflicts with the mass density observed in galaxies, which is considerably less
than \(\rho_c\). This discrepancy has led to an intensive search for signs of an
intergalactic gas, a search that has so far been quite unsuccessful.

Looking back in time, we find that any isotropic homogeneous universe governed
by Einstein's equations must have started with a singularity of infinite
density. Dating from this singularity, the age of the universe must be less than
\(H_0^{-1}\), and less than \(\tfrac23H_0^{-1}\) if \(q_0>\tfrac12\).
Radioactive dating and the theory of stellar evolution give uncertain ages,
ranging from \(7\times10^9\) to \(16\times10^9\) years, but it would be
difficult to admit an age much less than \(\tfrac23H_0^{-1}\).

The most prominent relic of the hot early universe is the \(2.7\,\mathrm{K}\)
microwave radiation background, predicted in 1950 and observed in 1965. The
weight of the data is so far consistent with the expectation that this radiation
has a Planck black-body spectrum and is perfectly isotropic. Knowing the present
radiation temperature, we can trace the thermal history of the universe back to
the first few minutes, and calculate the production of complex nuclei in the
primordial fireball. A fairly definite prediction emerges, that about 27\% of the
nucleons in the early universe should have been fused into
\({}^{4}\mathrm{He}\). This is in agreement with some measurements of the
present cosmic helium abundance, but in disagreement with others. Another relic
of the early universe is our present cosmic morphology: Stars form galaxies,
galaxies form clusters, and clusters form a more or less homogeneous gas. Our
present theoretical understanding of how this structure evolved is not in good
shape, but it is clear that the radiation background played an important role.
We can also speculate on the first few seconds of cosmic history, when the
temperature was high enough to produce mesons, baryons, and antibaryons in
large numbers; so far there does not seem to be any way to check the results of
these speculations.

The preceding summary describes what may be called the ``standard model'' of
the universe, based on the Cosmological Principle and Einstein's field equations.
One other ``standard'' assumption, which plays an important role in
Sections~\ref{sec:15.7}--\ref{sec:15.11}, is that the distant galaxies are, like
our own, composed of baryons rather than antibaryons. It has often been
suggested that since baryon number is, like charge, exactly conserved, the
universe ought to contain equal numbers of baryons and antibaryons as well as
positive and negative charges. However, it should be kept in mind that baryon
number is really unlike charge: There are long-range forces associated with
charge but, as far as we know, not with baryon number. Indeed, in a finite
universe the total charge must be zero, as can be immediately seen by integrating
the Maxwell equation \(\boldsymbol{\nabla}\cdot\mathbf{E}=\rho_{\mathrm{e}}\) over
the volume of the universe; no such conclusion can be derived for baryon
number. In any case, even if the net baryon number of the universe is zero,
baryons and antibaryons must somehow have become separated at some time in the
past, and most of the considerations of this chapter are applicable to the
evolution of the universe after that time.

Of course, the standard model may be partly or wholly wrong. However, its
importance lies not in its certain truth but in the common meeting ground that it
provides for an enormous variety of observational data. By discussing these data
in the context of a standard cosmological model, we can begin to appreciate their
cosmological relevance, whatever model ultimately proves correct. Some other
possible models are discussed in the next chapter.
